\chapter{Research Gap Identification}
\label{Research_Gap_Identification}

The literature reviewed in Chapter 4 identifies five gaps in the treatment of progress, visualisation and assistance in current educational platforms. Each gap is stated below together with the response adopted in the proposed system.

\section{Absence of Visual Progress Mapping}

Moodle, Canvas and Google Classroom represent a course as a list of assignments with associated scores. The representation conveys the value of a score but not the structure within which the work was performed, and a student cannot determine from it which topics were omitted or how the topics relate to one another.

In response, The list is replaced by an interactive map. Topics are represented as nodes, the syllabus supplies the edges, and the student may view the structure as a whole or examine a single branch. Three states are represented on the map, namely not started, in progress and completed, while the learning-path view derives a fourth state for topics that have been released and not opened. Both renderings are derived from the same records.

\section{Absence of Dual Progress Tracking}

Existing instruments record either the material covered by the teacher or the scores obtained by the student. Few present the two together, and the comparison a teacher requires, namely the position of the class relative to the position of the syllabus, is therefore left to recollection. An examination supplies that comparison only when the course is complete.

In response, Two records are maintained separately and compared when read. The first records the teacher's delivery; the second records each student's position. The comparison yields the topics released but not begun by a given student, together with an engagement classification for the class.

\section{Limited Application of Gamification}

Gamification has empirical support, but the systems that apply it effectively are specialised applications addressing a single subject and a single audience. General-purpose platforms have adopted badges and leaderboards instead of progression structures. A badge constitutes a reward; a skill tree is a representation of the course, and it is the representation that informs the student's decisions.

In response, The mindmap already has the form of a skill tree, so the progression mechanic requires no additional structure. Nodes move from not started to in progress and to completed as work is performed, and the learning-path view derives experience points, a level for each hundred points, and badges at the first, fifth, tenth and twenty-fifth completed topic.

\section{Absence of Curriculum Scope Control}

A general-purpose assistant answers from the corpus on which it was trained, which is not the corpus of a course. Such a system may be correct about the subject and incorrect about the syllabus within the same response, and the student has no means of distinguishing the two.

In response, Retrieval is restricted to one subject at query time. Passages are filtered by subject identifier, so the model can only be supplied with documents associated with that course, and a question whose terms match no indexed material receives an explicit statement that the material does not address it. The retrieval and generation path is implemented and tested; mounting the interface is the remaining step.

\section{Limited Multi-Tenant Support}

Most platforms assume a single institution. Serving several institutions from one installation requires either separate deployments or a shared database with a tenant column that is not enforced. The second arrangement is the more hazardous, since it resembles multi-tenancy until a query omits its tenant predicate.

In response, Row-Level Security is enabled on every table and tenant-scoped policies govern access through the browser client, where the role holds read access only. Application queries are issued through a dedicated non-owner role and carry an explicit institution predicate, which is the property that requires review. A test asserts that the subjects of two institutions partition the table exactly, that writes to curriculum tables are refused, and that anonymous reads are refused.

\section{Summary of Research Gaps}

\begin{table}[H]
\centering
\caption{Summary of Research Gaps and Proposed Solutions}
\label{tab:research_gaps}
\resizebox{\textwidth}{!}{
\begin{tabular}{|c|p{5cm}|p{5cm}|}
\hline
\thead{Gap \#} & \thead{Identified Gap} & \thead{Proposed Solution} \\ \hline
1 & No visual progress mapping in LMS & Interactive mindmap with node states \\ \hline
2 & No dual progress tracking & Teacher vs Student timeline comparison \\ \hline
3 & Limited gamification in LMS & Skill-tree style mindmap with XP and badges \\ \hline
4 & AI without curriculum scope & RAG scoped to teacher-approved materials \\ \hline
5 & No multi-tenant support & Row-Level Security for data isolation \\ \hline
\end{tabular}}
\end{table}
